% ══════════════════════════════════════════════════════════════════════════════
%  CHAPTER 1: FUNDAMENTALS
% ══════════════════════════════════════════════════════════════════════════════

\chapter{Fundamentals (lecture 1)}

% ──────────────────────────────────────────────────────────────────────────────
\section{Natural Units}

\begin{intuition}
In SI units, we treat time, length, mass, and energy as distinct physical quantities with different dimensions and units.
\smallskip\noindent
To make calculations more convenient, we introduce natural units to connect these quantities via fundamental constants,
so that we can express everything in terms of a single unit and same type of dimension (i.e.\ energy).
\end{intuition}

\begin{theory}
In natural units, we set $c = \hbar = k_B = 1$.

When $c=1$, mass and energy are measured in the same units and dimension in terms of $E^1$ by $ E = mc^2 = m $.
\end{theory}

\begin{example}{Dimension of Gravitational Constant}{}
What is the dimension of the gravitational constant $G$ in natural units?

\bigskip\textbf{Solution.}\quad
We know that $F = G \frac{m_1 m_2}{r^2}$ and $F = ma$. \\
We get dimension of $a$ is $E^1$ by $a = v/t = x/t^2 = \frac{E^{-1}}{E^{-1}\cdot E^{-1}} = E^1$. \\
Then we get dimension of $F$ is $E^2$ by $F = ma = E^1 \cdot E^1 = E^2$. \\
Finally, we get dimension of $G$ is $E^{-2}$ by $G = \frac{F r^2}{m_1 m_2} = E^2 \cdot E^{-2} / (E^1 \cdot E^1) = E^{-2}$.

\end{example}


In general, we define the dimension of a quantity directly in terms of n if its dimension is $E^n$ for simplicity.
Below are some useful conversion factors between SI and natural units:
\begin{center}
\begin{tabular}{@{}lll@{\quad}r@{}}
	\toprule
	\textbf{Quantity} & \textbf{Related equation} & \textbf{SI dimension} & \textbf{Natural unit} \\
\midrule
Energy & $E = mc^2$ & $\mathrm{kg\,m^2\,s^{-2}}$ & $1$ \\
Mass & $m = E/c^2$ & $\mathrm{kg}$ & $1$ \\
Frequency & $\nu = E/h$ & $\mathrm{Hz}$ & $1$ \\
Time & $t = 1/\nu$ & $\mathrm{s}$ & $-1$ \\
Space & $x = ct$ & $\mathrm{m}$ & $-1$ \\
Temperature & $T = E/k_B$ & $\mathrm{K}$ & $1$ \\
\midrule
Velocity & $v = x/t$ & $\mathrm{m\,s^{-1}}$ & $0$ \\
Momentum & $p = E/c$ & $\mathrm{kg\,m\,s^{-1}}$ & $1$ \\
Angular Velocity & $\omega = rv$ & $\mathrm{rad\,s^{-1}}$ & $-1$ \\
Angular Momentum & $L = \frac{1}{2}mv \omega$ & $\mathrm{kg\,m^2\,s^{-1}}$ & $0$ \\
Luminosity & $\mathcal{L} = E/t$ & $\mathrm{W}$ & $2$ \\
Force & $F = p/t$ & $\mathrm{N}$ & $2$ \\
Pressure & $P = F/A$ & $\mathrm{Pa}$ & $2$ \\
Mass Density & $\rho = m/V$ & $\mathrm{kg\,m^{-3}}$ & $4$ \\
Energy Density & $\rho = E/V$ & $\mathrm{J\,m^{-3}}$ & $4$ \\
\midrule
Gravitational Constant & $G = Fr^2/m_1m_2$ & $\mathrm{m^3\,kg^{-1}\,s^{-2}}$ & $-2$ \\
Redshift & $z = \lambda_0/\lambda_e - 1$ & --- & $0$ \\
Cosmological Scale Factor & $a(t) = 1/(1+z)$ & --- & $0$ \\
Hubble Parameter & $H = \dot{a}/a$ & $\mathrm{s^{-1}}$ & $1$ \\
\bottomrule
\end{tabular}
\end{center}

\bigskip\noindent
\bigskip\noindent
\begin{remark}
  In the Heavide-Lorentz system of natural units, we also set $\epsilon_0 = \mu_0 = 1$.\\
  Which one can hence get the dimension of charge $q$ is 0.
\end{remark} 



% ──────────────────────────────────────────────────────────────────────────────
\newpage
\section{4-Vectors and Tensors}

\subsection{4-Vectors}
\begin{foundation}
  In special relativity, we treat the time and space coordinates on equal footing as a 4-vector $x^\mu = (t, x, y, z) = (x^0, x^1, x^2, x^3)$.
  or in other words, we used the 4-vector to combine the time and space coordinates together.
\end{foundation}

\begin{theory}
  By using the 4-vector, we can write the Lorentz transformation in-terms of the Lorentz transformation matrix and the 4-vector: \\
  \smallskip\noindent
  WLOG, we can consider a boost in the $x$-direction with velocity $v$,
  \[
    \begin{pmatrix}t' \\ x' \\ y' \\ z'
    \end{pmatrix}
    = \begin{pmatrix}
      \gamma & -\gamma \beta & 0 & 0 \\
      -\gamma \beta & \gamma & 0 & 0 \\
      0 & 0 & 1 & 0 \\
      0 & 0 & 0 & 1
    \end{pmatrix}
    \begin{pmatrix}t \\ x \\ y \\ z
    \end{pmatrix},
  \]
  where $\gamma = 1/\sqrt{1-v^2}$ is the Lorentz factor, and $\beta = v/c$ is the velocity in units of $c$ ($\beta = v$ in natural units).
\end{theory}

% ──────────────────────────────────────────────────────────────────────────────
\bigskip\noindent
\subsection{Einstein Summation Convention}
\begin{intuition}
  Each time when we are doing the 4-vectors operations like the Lorentz transformation, we are doing the summation of the products of the components of the 4-vectors and the Lorentz transformation matrix. \\
  \smallskip\noindent
  It can be very tedious to write down the summation every time, so we want a more compact notation to represent the summation.
\end{intuition}
\begin{theory}
  Using the example of the Lorentz transformation above, if we want to do the transformation in traditional way, we have to write down the summation explicitly:
  \[
    x^{\mu'} = \Sigma_{\nu=0}^3 \Lambda^\mu{}_\nu x^\nu.
  \]
  \smallskip\noindent
  By using the Einstein summation convention, we can simply write down the transformation as:
  \[
    x^{\mu'} = \Lambda^\mu{}_\nu x^\nu.
  \]
\end{theory}


% ──────────────────────────────────────────────────────────────────────────────
\bigskip\noindent
\subsection{Metric Tensor}
\begin{intuition}
  In a non-Euclidean space, such as the Minkowski spacetime in special relativity, the distance between two points is not simply what we seen in our eyes. \\
  The spacetime interval between two events(or points) is no longer measured by the Pythagorean theorem, so we need to use another ruler to measure it. \\
  \\ 
  In general, we used a metric tensor as the ruler to measure the distance of 2 point in a specific spacetime.
\end{intuition}

\begin{theory}
  In traditional Euclidean space, the metric tensor is simply the identity matrix, so the distance between two points is given by the Pythagorean theorem:
  \[
    ds^2 = \begin{pmatrix}dx & dy & dz\end{pmatrix}
    \begin{pmatrix}1 & 0 & 0 \\
            0 & 1 & 0 \\
            0 & 0 & 1\end{pmatrix}
    \begin{pmatrix}dx \\ dy \\ dz\end{pmatrix}
    = dx^2 + dy^2 + dz^2.
  \]
  \smallskip\noindent
  In Minkowski spacetime, the metric tensor is given by:
  \[
    g_{\mu\nu} = \begin{pmatrix}-1 & 0 & 0 & 0 \\
            0 & 1 & 0 & 0 \\
            0 & 0 & 1 & 0 \\
            0 & 0 & 0 & 1\end{pmatrix},
  \]
  so the spacetime interval is given by:
  \[
    ds^2 = g_{\mu\nu} dx^\mu dx^\nu = -dt^2 + dx^2 + dy^2 + dz^2.
  \]
\end{theory}

  or in other words, the spacetime interval is calculated by the dot product of the 4-vector with itself by using the metric tensor as the dot product operator:
\begin{theory}
  \[
    A \cdot B = g_{\mu\nu} A^\mu B^\nu = \begin{pmatrix} A^0 & A^1 & A^2 & A^3 \end{pmatrix}
    \begin{pmatrix}
      g_{00} & g_{01} & g_{02} & g_{03} \\
      g_{10} & g_{11} & g_{12} & g_{13} \\
      g_{20} & g_{21} & g_{22} & g_{23} \\
      g_{30} & g_{31} & g_{32} & g_{33}
    \end{pmatrix}
    \begin{pmatrix} B^0 \\ B^1 \\ B^2 \\ B^3 \end{pmatrix}.
  \]
\end{theory}


% ──────────────────────────────────────────────────────────────────────────────
\newpage
\subsection{Tensors}
\begin{intuition}
  Now you may wonder, why sometimes the index is up, and sometimes it is down? Is there any difference between them? \\
  \bigskip
  The answer is yes, they are different. \\
  The upper index represents a \textbf{contravariant} component, while the lower index represents a \textbf{covariant} component. \\
  By the combination of the \textbf{contravariant} and \textbf{convariant} components, we can construct a more general mathematical object called a \textbf{tensor}.
\end{intuition}

\bigskip 
\begin{foundation}
  In tensor calculus and differential geometry, the position of an index (upper or lower) indicates how the component of the tensor transforms when you change the coordinate system.

  \bigskip\noindent
  \textbf{Direct Mapping:}
  \begin{itemize}
    \item \textbf{Upper Index} ($\mu$ in $V^\mu$): Represents a \textbf{Contravariant} component.
    \item \textbf{Lower Index} ($\mu$ in $V_\mu$): Represents a \textbf{Covariant} component.
  \end{itemize}

  \bigskip\noindent
  \noindent\textbf{Transformation Rules:}

  The terms \emph{contravariant} and \emph{covariant} describe how components change relative to \textbf{basis vectors} when switching coordinates (e.g., Cartesian to Polar, or meters to centimeters).

  \bigskip\noindent
  \noindent\textit{Contravariant (Upper Index):}
  \begin{itemize}
    \item \textbf{Symbol:} $V^\mu$
    \item \textbf{Transformation:} Transforms \textbf{inversely} (against) the basis vectors.
    \medskip
    \item \textbf{Law:} $\displaystyle V'^\mu = \frac{\partial x'^\mu}{\partial x^\nu} V^\nu$
    \medskip
    \item \textbf{Intuition:} If your ruler shrinks (smaller basis vectors), the numerical value must grow to describe the same physical length. Example: $1 \text{ m} = 100 \text{ cm}$.
  \end{itemize}

  \bigskip\noindent
  \noindent\textit{Covariant (Lower Index):}
  \begin{itemize}
    \item \textbf{Symbol:} $V_\mu$
    \item \textbf{Transformation:} Transforms \textbf{in the same way} (with) the basis vectors.
    \medskip
    \item \textbf{Law:} $\displaystyle V'_\mu = \frac{\partial x^\nu}{\partial x'^\mu} V_\nu$
    \medskip
    \item \textbf{Intuition:} Like a gradient (slope). If the ruler shrinks, the slope per unit distance also shrinks.
  \end{itemize}

  \bigskip\noindent
  \noindent\textbf{Geometric Interpretation:}
  \smallskip\noindent
  \begin{itemize}
    \item \textbf{Contravariant Vectors} ($V^\mu$): Standard ``arrows'' tangent to a curve, living in the \textbf{Tangent Space}.
    \medskip
    \item \textbf{Covariant Vectors} ($V_\mu$): Called \textbf{Dual Vectors} or \textbf{One-Forms}. They are functions that map vectors to scalars and live in the \textbf{Cotangent Space}. Geometrically visualized as stacked parallel surfaces pierced by vectors.
  \end{itemize}

  \bigskip\noindent
  \noindent\textbf{Einstein Summation Convention:} \\
  \smallskip\noindent
  The distinction is critical, repeated indices (one upper, one lower) indicate summation. Only these are coordinate-invariant scalars.
  \medskip
  \begin{itemize}
    \item \textbf{Valid:} $S = A^\mu B_\mu$ \quad $\Rightarrow$ scalar invariant
    \item \textbf{Invalid:} $A^\mu B^\mu$ or $A_\mu B_\mu$ \quad $\Rightarrow$ not coordinate-invariant (not summation)
  \end{itemize}

  \bigskip\noindent
  \noindent\textbf{Summary Table}
  \bigskip\noindent
  \begin{center}
  \begin{tabular}{@{}lll@{}}
    \toprule
    \textbf{Property} & \textbf{Contravariant} & \textbf{Covariant} \\
    \midrule
    Index Position & Upper ($^\mu$) & Lower ($_\mu$) \\
    Notation & $V^\mu$ & $V_\mu$ \\
    Also Called & Vector & Dual Vector / One-Form \\
    Transformation & Opposite to basis & Same as basis \\
    Example & Displacement, Velocity & Gradient \\
    Metric Operation & Lowered by $g_{\mu\nu}$ & Raised by $g^{\mu\nu}$ \\
    \bottomrule
  \end{tabular}
  \end{center}

\end{foundation}

\begin{theory}
  \bigskip\noindent
  \noindent\textbf{Raising and Lowering Indices:} \\
  \smallskip\noindent
  With the \textbf{Metric Tensor} ($g_{\mu\nu}$), you can convert between contravariant and covariant components:
  \medskip
  \begin{itemize}
    \item \textbf{Lowering:} $\quad V_\mu = g_{\mu\nu} V^\nu$ \quad (contravariant $\to$ covariant)
    \item \textbf{Raising:} $\quad V^\mu = g^{\mu\nu} V_\nu$ \quad (covariant $\to$ contravariant)
  \end{itemize}
\end{theory}

\newpage
And the contravariant(inverse) metric $g^{\mu\nu}$ is the inverse of the covariant metric $g_{\mu\nu}$, so we have:
\begin{theory}
  \[
    g^{\mu\nu} g_{\nu\sigma} = \delta^\mu{}_\sigma,
  \]
  where $\delta^\mu{}_\sigma$ is the Kronecker delta, which is 1 if $\mu = \sigma$ and 0 otherwise.
\end{theory}


As a results, we can further simplified the operations like the dot product and trace:
\begin{theory}
  \[
    A \cdot B = g_{\mu\nu} A^\mu B^\nu = A_\nu B^\nu = A^\mu B_\mu.
  \]
  \[
    T = g^{\mu\nu} T_{\mu\nu} = g_{\mu\nu} T^{\mu\nu} = T^\nu{}_\nu = T_\mu{}^\mu.
  \]
\end{theory}  

\bigskip\noindent
\begin{example}{Trace of Spherical Minkowski Metric}{}
  Calculate the trace of the spherical Minkowski metric $g_{\mu\nu} = \operatorname{diag}(-1, 1, r^2, r^2 \sin^2\theta)$.

  \bigskip\textbf{Solution.}\quad
  We need to first calculate the inverse metric $g^{\mu\nu}$, which is given by:
  \[
    g^{\mu\nu} = g_{\mu\nu}^{-1} = \operatorname{diag}(-1, 1, r^{-2}, r^{-2} \sin^{-2}\theta).
  \]
  Then we can calculate the trace:
  \[
    T = g^{\mu\nu} g_{\mu\nu} = (-1)(-1) + (1)(1) + (r^{-2})(r^2) + (r^{-2} \sin^{-2}\theta)(r^2 \sin^2\theta) = 4.
  \]
\end{example}

\medskip\noindent
\begin{importantnote}
  In flat Euclidean space with Cartesian coordinates, $g_{\mu\nu} = \delta_{\mu\nu}$ (identity), so $V^\mu = V_\mu$ numerically. 

  \bigskip\noindent
  However, for other spacetimes (e.g., Minkowski, curved), the metric is not the identity, and $V^\mu$ and $V_\mu$ can differ significantly. The distinction becomes crucial in General Relativity and curvilinear coordinates. To be safe, treat them as different objects and always check the metric when raising/lowering indices.
\end{importantnote}

\bigskip\noindent
\begin{extra}
  The tensor T is said to be of type (p,q) if it has:
  \begin{itemize}
    \item p: The number of contravariant indices(upper indices)
    \item q: The number of convariant indices(lower indices)
  \end{itemize}
  Tensor with different types cannot be added together or equal to each other, in such a meaning, you can think of the type of a tensor as its ``dimension'' in the space of tensors, and only tensors with the same type can be compared or combined.

  \bigskip\noindent
  \textbf{Tensor Rank:}  The total rank of a Tensor T is p + q (let said p + q = n), and we will said it is a rank-n-Tensor \\

  \smallskip\noindent
  Note that the rank of a tensor is not the same as linear algebra rank of a matrix. The rank of a tensor is the total number of indices, while the linear algebra rank is the maximum number of linearly independent rows or columns in a matrix. Thats mean a rank-2 tensor can have a linear algebra rank of 0, 1, or 2 depending on its components.

  \bigskip\noindent
  \begin{center}
    \begin{tabular}{@{}c|c|c|c|l@{}}
      \toprule
      \textbf{Tensor Type} & \textbf{p} & \textbf{q} & \textbf{Rank} & \textbf{Example} \\
      \midrule
      (0,0) & 0 & 0 & 0 & Scalar (e.g., Temperature) \\
      (1,0) & 1 & 0 & 1 & Contravariant Vector (e.g., Displacement $V^\mu$) \\
      (0,1) & 0 & 1 & 1 & Covariant Vector (e.g., Gradient $V_\mu$) \\
      (1,1) & 1 & 1 & 2 & Mixed Tensor (e.g., Linear Transformation $\Lambda^\mu{}_\nu$) \\
      (2,0) & 2 & 0 & 2 & Contravariant Tensor (e.g., Inverse Metric $g^{\mu \nu}$) \\
      (0,2) & 0 & 2 & 2 & Covariant Tensor (e.g., Metric Tensor $g_{\mu \nu}$) \\
      (1,3) & 1 & 3 & 4 & Riemann Curvature ($R^{\rho}_{\sigma \mu \nu}$) \\
      \bottomrule
    \end{tabular}
  \end{center}
\end{extra}

\clearpage